\documentclass[10pt]{beamer}
\usepackage[utf8]{inputenc}
\usepackage[T1,T2A]{fontenc}
\usepackage[russian]{babel}
\usepackage{color}
\usepackage{calc}
\usepackage{graphicx}
\usepackage{epstopdf}
\usepackage{hyperref}
\hypersetup{unicode,colorlinks}
\usepackage{csquotes}
\usepackage{upquote}
\usepackage{cprotect}
\usetheme[progressbar=head,numbering=fraction,block=fill]{metropolis}
\usepackage{dejavu}
\usepackage{etoolbox}
\usepackage{bm}
\usepackage[ND]{prftree}
\usepackage[tableaux]{prooftrees}
\usepackage{mathtools} % for bigtimes
\usepackage{diagbox}
\usepackage{verbatimbox}
\usepackage{fitch}
% fitch по умолчанию печатает эти правила как $\Rightarrow$I/E; в курсе импликация всюду $\to$.
\gappto\ndrules{\def\ii{\by{$\to$I}}\def\ie{\by{$\to$E}}}
%\usepackage[defaultsans]{droidsans}
%\usepackage[defaultmono]{droidmono}
%\makeatletter
%\input droid/t1fdm.fd
%\makeatother

\makeatletter
\DeclareRobustCommand*{\bfseries}{%
    \not@math@alphabet\bfseries\mathbf
    \fontseries\bfdefault\selectfont
    \boldmath
}
\makeatother

\makeatletter
\chardef\straightquote@code=\catcode`'
\chardef\backquote@code=\catcode``
\catcode`'=\active \catcode``=\active
\patchcmd{\@noligs}
{\textasciigrave}
{\fixedtextasciigrave}
{}{}
\newcommand{\fixedtextasciigrave}{%
    \makebox[.5em]{\fontencoding{TS1}\fontfamily{fvs}\selectfont\textasciigrave}% Vera Sans
}
\catcode\lq\'=\straightquote@code
\catcode\lq\`=\backquote@code
\makeatother

\setbeamercovered{again covered={\opaqueness<1->{25}}}

\forestset{close with=\ensuremath{\bigtimes}}

\newcommand{\customframefont}[1]{
    \setbeamertemplate{itemize/enumerate body begin}{#1}
    \setbeamertemplate{itemize/enumerate subbody begin}{#1}
}

\NewEnviron{framefont}[1]{
    \customframefont{#1} % for itemize/enumerate
    {#1 % For the text outside itemize/enumerate
        \BODY
    }
    \customframefont{\normalsize}
}


\makeatletter
\DeclareRobustCommand{\rvdots}{%
    \vphantom{\int\limits^x}\smash[t]{\vdots}
}
\makeatother

% Короткое вертикальное многоточие: размеры в em, поэтому оно уменьшается вместе со шрифтом
% (у \vdots они фиксированы в pt).
\DeclareRobustCommand{\ndvdots}{%
    \mathord{\vbox{\baselineskip=.3em\lineskiplimit=0pt\kern.15em\hbox{.}\hbox{.}\hbox{.}}}}
% Подвывод в схеме правила, нарисованный как во fitch: вертикальная черта и флажок под допущением.
% Выравнивание по нижней строке, чтобы в \lor E заключения подвыводов стояли на одной линии с A \lor B.
\newcommand{\ndsubproof}[2]{{%
    \setlength{\arrayrulewidth}{.2mm}%
    \begin{array}[b]{@{}|@{\,}l@{\,}}#1\\\hline\ndvdots\\#2\end{array}}}

\newcommand{\N}{\mathbb{N}}
\newcommand{\Z}{\mathbb{Z}}
\newcommand{\Q}{\mathbb{Q}}
\newcommand{\R}{\mathbb{R}}
\newcommand{\pow}[1]{\mathcal{P}({#1})}
\newcommand{\continuum}{\mathfrak{c}}

% \reveal<спец>{...}: показать содержимое с указанного слайда, до этого оставив пустое место той же ширины.
% В отличие от \pause и \uncover не оставляет невидимых PGF-специалов; внутри выключной формулы \pause
% оставляет их в пустом абзаце после формулы, и TeX набирает лишнюю пустую строку.
% Следующий \item после такой формулы получает <.(1)-> (тот же слайд, что и \reveal<+(1)->)
% или <+(1)-> (следующий слайд); в beamer «.» "--- это значение счётчика минус один.
\newcommand<>{\reveal}[1]{\alt#2{#1}{\phantom{#1}}}

% Многобуквенные имена в формулах: курсив с нормальным кернингом (а не произведение букв).
\newcommand{\ident}[1]{\mathit{#1}}
\newcommand{\Var}{\ident{Var}}
\newcommand{\Prop}{\ident{Prop}}
\newcommand{\Const}{\ident{Const}}
\newcommand{\Fun}{\ident{Fun}}
\newcommand{\Pred}{\ident{Pred}}
\newcommand{\Term}{\ident{Term}}
\newcommand{\Frm}{\ident{Form}} % \Form занят hyperref
\newcommand{\arity}{\ident{arity}}
\newcommand{\Th}{\ident{Th}}
\newcommand{\Dom}{\ident{Dom}}
\newcommand{\Ran}{\ident{Ran}}
\newcommand{\id}{\ident{id}}
\newcommand{\Ord}{\ident{Ord}}
\newcommand{\RAA}{\ident{RAA}}
% Названия теорий и аксиом: прямой шрифт.
\newcommand{\ZF}{\mathrm{ZF}}
\newcommand{\ZFC}{\mathrm{ZFC}}
\newcommand{\AC}{\mathrm{AC}}
\newcommand{\DC}{\mathrm{DC}}
\newcommand{\CH}{\mathrm{CH}}
\newcommand{\GCH}{\mathrm{GCH}}

\author{Алексей Романов}
\subtitle{Математическая логика и теория алгоритмов}
%\logo{}
\institute{МИЭТ}
\subject{Математическая логика и теория алгоритмов}
%\setbeamercovered{transparent}
%\setbeamertemplate{navigation symbols}{}


\title{Множества и функции}
\date{\today}

\begin{document}

\begin{frame}[plain]
    \maketitle
\end{frame}

\begin{frame}
    \frametitle{Задаёт ли запись функцию?}
    \begin{itemize}
        \item Указаны домен (область определения) $A$, кодомен (область прибытия) $B$ и правило $f(x)=\ldots$
        \item Для каждого $x\in A$ должно быть ровно одно значение в $B$.
        \item<2-> Проверяем:
        \begin{enumerate}
            \item<2-> значение существует для каждого $x\in A$;
            \item<3-> оно принадлежит $B$;
            \item<4-> оно однозначно определено.
        \end{enumerate}
        \item<5-> Если правило задано формулой логики предикатов $E(x,y)$, требуется $\forall x:A~\exists!y:B~E(x,y)$.
        \item<6-> Если запись не задаёт функцию $A\to B$, вопрос о её инъективности и сюръективности не ставится.
        \item Здесь $\N=\{0,1,2,\ldots\}$.
    \end{itemize}
\end{frame}

\begin{frame}
    \frametitle{Существование и допустимость значения}
    \begin{itemize}
        \item Задаёт ли $f:\R\to\R$, $f(x)=1/x$ функцию?
        \item<2-> Нет: при $x=0$ значение не определено.
        \item<3-> А $f:\R\setminus\{0\}\to\R$, $f(x)=1/x$?
        \item<4-> Да: для каждого допустимого $x$ значение единственно и действительно.
        \item<5-> Задаёт ли $g:\N\to\N$, $g(n)=n/2$ функцию?
        \item<6-> Нет: $g(1)=1/2\notin\N$.
        \item<7-> Но та же формула задаёт функцию $\N\to\Q$.
    \end{itemize}
\end{frame}

\begin{frame}
    \frametitle{Уравнение не всегда определяет функцию}
    \begin{itemize}
        \item Условие $f(x)^2=x$ должно определить $f:A\to B$.
        \item Это задание правила формулой $E(x,y)\equiv y^2=x$.
        \item Для каких домена и кодомена условие однозначно задаёт функцию (то есть \pause $\forall x:A~\exists!y:B~y^2=x$)?
        \item<3-> $A=B=\R$: \uncover<4->{при $x=-1$ подходящего значения нет.}
        \item<5-> $A=[0,+\infty)$, $B=\R$: \uncover<6->{при $x=1$ подходят $1$ и $-1$.}
        \begin{itemize}
            \item<7-> Само условие не выбирает одно из этих значений.
        \end{itemize}
        \item<8-> $A=B=[0,+\infty)$: \uncover<9->{для каждого $x$ есть ровно одно подходящее значение.}
        \begin{itemize}
            \item<10-> Получаем функцию $f(x)=\sqrt{x}$.
        \end{itemize}
    \end{itemize}
\end{frame}

\begin{frame}
    \frametitle{Один аргумент, разные записи}
    \begin{itemize}
        \item Попробуем задать $f:\Q\to\Z$ правилом $f(n/m)=nm$, где $n\in\Z$, $m\in\Z\setminus\{0\}$.
        \item Определено ли значение однозначно?
        \item<2-> $1/2=2/4$ "--- один и тот же рациональный аргумент.
        \item<3-> Но правило даёт $f(1/2)=1\cdot2=2$ и $f(2/4)=2\cdot4=8$.
        \item<4-> Значит, без дополнительных условий это не определение функции.
        \item<5-> Когда аргумент имеет несколько представлений, результат не должен зависеть от выбранного представления.
        \item<6-> Можно исправить правило: $f(n/m)=\dfrac{nm}{\operatorname{\text{НОД}}(|n|,|m|)^2}$. Проверьте, что это функция $\Q\to\Z$: результат целый и не зависит от представления дроби.
    \end{itemize}
\end{frame}

\begin{frame}
    \frametitle{Правило на несократимых дробях}
    \begin{itemize}
        \item Теперь в $f(n/m)=nm$ требуем $m>0$ и $\operatorname{\text{НОД}}(|n|,m)=1$.
        \item Получили ли функцию $\Q\to\Z$?
        \item<2-> Да: у каждого рационального числа единственное такое представление. Для $0$ это $0/1$.
        \item<3-> Является ли эта функция вложением (инъекцией)?
        \item<4-> Нет: $f(1/2)=2=f(2)$, хотя $1/2\neq2$.
        \item<5-> Является ли она наложением (сюръекцией)? Для каждого ли $k\in\Z$ есть прообраз?
        \item<6-> $k=k/1$ "--- несократимая дробь, и $f(k)=k\cdot1=k$.
        \item<7-> Это наложение, но не вложение и не биекция.
    \end{itemize}
\end{frame}

\begin{frame}
    \frametitle{Как проверять инъективность и сюръективность}
    \begin{itemize}
        \item Пусть уже установлено, что $f:A\to B$ "--- функция.
        \item<2-> \emph{Инъективность:}
        \uncover<3->{$\forall x,y:A~(f(x)=f(y)\to x=y)$.
        Для произвольных $a,b\in A$ из $f(a)=f(b)$ выводим $a=b$.}
        \item<4-> \emph{Неинъективность:}
        \uncover<5->{$\exists x,y:A~(x\neq y\land f(x)=f(y))$.
        Находим $a,b\in A$ такие, что $a\neq b$ и $f(a)=f(b)$.}
        \item<6-> \emph{Сюръективность:}
        \uncover<7->{$\forall y:B~\exists x:A~f(x)=y$.
        Для произвольного $b\in B$ находим $a\in A$ с $f(a)=b$. Обычно выражаем $a$ термом, зависящим от $b$.}
        \item<8-> \emph{Несюръективность:}
        \uncover<9->{$\exists y:B~\forall x:A~f(x)\neq y$.
        Находим $b\in B$, для которого $f(a)\neq b$ при произвольном $a\in A$.}
        \item<10-> \emph{Биекция:} доказываем оба свойства. Позже увидим способ через обратную функцию.
    \end{itemize}
\end{frame}

\begin{frame}
    \frametitle{Одна формула, разные кодомены}
    \begin{itemize}
        \item $f:\R\setminus\{0\}\to\R$, $f(x)=1/x$.
        \item Проверим инъективность и сюръективность.
        \item<2-> Рассмотрим произвольные $a,b\in\R\setminus\{0\}$ такие, что $f(a)=f(b)$. Тогда $1/a=1/b$.
        \item<3-> Умножаем на $ab$: $b=a$. Это вложение.
        \item<4-> Сюръективности нет: $0$ не имеет прообраза.
        \item<5-> Заменим кодомен на $\R\setminus\{0\}$.
        \item<6-> Для произвольного $b\in\R\setminus\{0\}$ берём $a=1/b\in\R\setminus\{0\}$. Тогда $f(a)=b$.
        \item<7-> Теперь это биекция. Правило осталось тем же!
    \end{itemize}
\end{frame}

\begin{frame}
    \frametitle{Удвоение натурального числа}
    \begin{itemize}
        \item $f:\N\to\N$, $f(n)=2n$.
        \item Проверим инъективность и сюръективность.
        \item<2-> Сначала проверяем функцию: $2n\in\N$ для каждого $n\in\N$, значение единственно.
        \item<3-> Рассмотрим произвольные $a,b\in\N$ такие, что $f(a)=f(b)$. Тогда $2a=2b$.
        \item<4-> Значит, $a=b$: инъективность доказана.
        \item<5-> Каждое ли число из $\N$ имеет прообраз?
        \item<6-> $1$: уравнение $2n=1$ не имеет решений в $\N$.
        \item<7-> Поэтому это не наложение и не биекция.
    \end{itemize}
\end{frame}

\begin{frame}
    \frametitle{Удвоение рационального числа}
    \begin{itemize}
        \item Теперь $f:\Q\to\Q$, $f(x)=2x$.
        \item Что изменилось по сравнению с $\N\to\N$?
        \item<2-> Это функция: для каждого $x\in\Q$ значение $2x$ единственно и принадлежит $\Q$.
        \item<3-> Для произвольных $a,b\in\Q$ из $f(a)=f(b)$ получаем $2a=2b$, откуда $a=b$.
        \item<4-> Для произвольного $b\in\Q$ нужно найти $a\in\Q$ с $2a=b$.
        \item<5-> Берём $a=b/2\in\Q$. Тогда $f(a)=2(b/2)=b$.
        \item<6-> Значит, это наложение, следовательно, биекция.
        \item<7-> Для $\N\to\N$ этот шаг не работал: например, $1/2\notin\N$. Для $\Q\to\Q$ деление на $2$ не выводит из домена.
    \end{itemize}
\end{frame}

\begin{frame}
    \frametitle{Деление с округлением вниз}
    \begin{itemize}
        \item $g:\N\to\N$, $g(2k)=g(2k+1)=k$ для $k\in\N$.
        \item Эквивалентно, $g(n)=\lfloor n/2\rfloor$.
        \item Проверим, задаёт ли это правило функцию $\N\to\N$.
        \item<2-> Каждое натуральное число единственным образом равно $2k$ или $2k+1$. Значение $g$ определено и лежит в $\N$.
        \item<3-> Проверим инъективность и сюръективность.
        \item<4-> Для произвольного $b\in\N$ нужно найти прообраз.
        \item<5-> Берём $a=2b\in\N$: $g(a)=g(2b)=b$. Это наложение.
        \item<6-> Но $g(0)=g(1)=0$, хотя $0\neq1$.
        \item<7-> Значит, это не вложение и не биекция.
    \end{itemize}
\end{frame}

\begin{frame}
    \frametitle{Остаток от деления}
    \begin{itemize}
        \item $h:\N\to\N$, $h(n)=n\bmod2$.
        \item Определите, является ли $h$ вложением, наложением, биекцией.
        \item<2-> Это функция: каждый аргумент имеет единственный остаток $0$ или $1$, принадлежащий $\N$.
        \item<3-> $h(0)=h(2)=0$: не вложение.
        \item<4-> $h(n)\in\{0,1\}$ для любого $n$. Например, $2$ не имеет прообраза: не наложение.
        \item<5-> Значит, не биекция.
        \item<6-> А если кодомен заменить на $\{0,1\}$?
        \item<7-> Получим наложение: $h(0)=0$, $h(1)=1$. Инъективности по-прежнему нет.
    \end{itemize}
\end{frame}

\begin{frame}
    \frametitle{Проекция на первый компонент}
    \begin{itemize}
        \item Пусть $A,B\neq\varnothing$. Рассмотрим $\pi:A\times B\to A$, $\pi(a,b)=a$.
        \item Для каких $A,B$ это вложение, наложение, биекция?
        \item<2-> Это функция: первый компонент каждой пары однозначно определён и принадлежит $A$.
        \item<3-> Сюръективность: фиксируем $b_0\in B$. Для каждого $a\in A$ имеем $\pi(a,b_0)=a$.
        \item<4-> Если в $B$ есть $b_1\neq b_2$, то $\pi(a,b_1)=\pi(a,b_2)$ для любого $a\in A$. Инъективности нет.
        \item<5-> Если $B=\{b_0\}$, то для произвольных $a,a'\in A$ из $\pi(a,b_0)=\pi(a',b_0)$ следует равенство пар.
        \item<6-> При непустых $A,B$ проекция биективна тогда и только тогда, когда $B$ одноэлементно.
    \end{itemize}
\end{frame}

\begin{frame}
    \frametitle{Проекция: пустые множества}
    \begin{itemize}
        \item Что изменится для $\pi:A\times B\to A$, если одно из множеств пусто?
        \item<2-> Если $A=\varnothing$, то $A\times B=\varnothing$ для любого $B$.
        \item<3-> Получаем пустую функцию $\varnothing\to\varnothing$. Она биективна.
        \item<4-> Если $A\neq\varnothing$, но $B=\varnothing$, то $A\times B=\varnothing$.
        \item<5-> Пустая функция $\varnothing\to A$ является вложением: различных аргументов нет.
        \item<6-> Она не является наложением: ни один элемент непустого $A$ не имеет прообраза.
        \item<7-> В задачах о произвольных множествах пустые случаи нужно проверять отдельно.
    \end{itemize}
\end{frame}

\begin{frame}
    \frametitle{Многочлен пятой степени}
    \begin{itemize}
        \item $f:\R\to\R$, $f(x)=x^5+x$.
        \item Проверим инъективность и сюръективность.
        \item<2-> $f'(x)=5x^4+1>0$ при любом $x\in\R$. Поэтому $f$ строго возрастает.
        \item<3-> Разные аргументы дают разные значения: $f$ "--- вложение. Одного неубывания для этого недостаточно.
        \item<4-> Многочлен $f$ непрерывен, и $f(x)\to\pm\infty$ при $x\to\pm\infty$.
        \item<5-> Для произвольного $b\in\R$ существуют $a<c$, для которых $f(a)<b<f(c)$.
        \item<6-> По теореме о промежуточных значениях существует $x\in(a,c)$ такой, что $f(x)=b$.
        \item<7-> Значит, $f$ "--- наложение и биекция.
        \item<7-> Явная формула прообраза для этого доказательства не нужна.
    \end{itemize}
\end{frame}

\begin{frame}
    \frametitle{Тот же многочлен на $\Q$}
    \begin{itemize}
        \item Теперь $f:\Q\to\Q$, $f(x)=x^5+x$.
        \item Проверим инъективность и сюръективность.
        \item<2-> Значения рациональны. Строгое возрастание сохраняется, поэтому $f$ "--- вложение.
        \item<3-> Есть ли у $1\in\Q$ рациональный прообраз?
        \item<4-> Допустим, $x=n/m$, где $m>0$ и $\operatorname{\text{НОД}}(|n|,m)=1$. Тогда $n^5+nm^4=m^5$.
        \item<5-> Отсюда $n^5=m^4(m-n)$, значит, $m\mid n^5$.
        \item<6-> Несократимость даёт $m=1$, так что $x=n\in\Z$.
        \item<7-> Но $n^5+n$ чётно при любом целом $n$, поэтому не равно $1$. Противоречие!
        \item<8-> Значит, $f$ не является наложением и биекцией.
        \item<9-> Теорема о промежуточных значениях даёт прообраз в $\R$, но не обязательно в $\Q$.
    \end{itemize}
\end{frame}

\begin{frame}
    \frametitle{Обратная функция}
    \begin{itemize}
        \item Пусть $f:A\to B$. Что значит, что $g$ "--- обратная к $f$?
        \item<2-> Это функция $g:B\to A$, для которой
        \begin{itemize}
            \item<3-> $g\circ f=\id_A$, то есть $\forall x:A~g(f(x))=x$;
            \item<4-> $f\circ g=\id_B$, то есть $\forall y:B~f(g(y))=y$.
        \end{itemize}
        \item<5-> Обозначение: $g=f^{-1}$. Проверять нужно обе композиции.
        \item<6-> Обратная существует ровно у биекций: каждому $y\in B$ сопоставляем его единственный прообраз.
        \item<7-> Не путать с прообразом множества $f^{-1}[C]$: для него обратная функция не требуется.
    \end{itemize}
\end{frame}

\begin{frame}
    \frametitle{Поиск прообразов и обратная функция}
    \begin{itemize}
        \item Пусть $f:A\to B$. Обычный способ доказать сюръективность "--- для произвольного $y\in B$ найти прообраз $x=g(y)\in A$.
        \item<2-> Так строим функцию $g:B\to A$, для которой $f(g(y))=y$ при каждом $y\in B$.
        \item<3-> То есть $f\circ g=\id_B$: $g$ "--- \emph{правая обратная} к $f$. Этого достаточно для сюръективности $f$.
        \item<4-> Для биективности проверим, что та же $g$ является и \emph{левой обратной}: $g(f(x))=x$ при каждом $x\in A$.
        \item<5-> Тогда $g\circ f=\id_A$, значит, $g=f^{-1}$ и $f$ "--- биекция.
        \item<6-> Для произвольных $a,b\in A$ из $f(a)=f(b)$ получаем $g(f(a))=g(f(b))$, то есть $a=b$: $f$ инъективна.
    \end{itemize}
\end{frame}

\begin{frame}
    \frametitle{Дополнение подмножества}
    \begin{itemize}
        \item Для любого множества $A$ рассмотрим $c:\pow{A}\to\pow{A}$, $c(X)=A\setminus X$.
        \item Проверим инъективность и сюръективность. Попробуем найти обратную функцию.
        \item<2-> Это функция: при $X\subseteq A$ имеем $A\setminus X\subseteq A$.
        \item<3-> Как из $A\setminus X$ восстановить $X$?
        \item<4-> Ещё раз взять дополнение: $d=(Y\mapsto A\setminus Y)$.
        \item<5-> Для произвольного $X\subseteq A$: $d(c(X))=A\setminus(A\setminus X)=X$.
        \item<6-> Для произвольного $Y\subseteq A$: $c(d(Y))=A\setminus(A\setminus Y)=Y$.
        \item<7-> Значит, $c$ "--- биекция, причём $c^{-1}=c$. Это верно и при $A=\varnothing$.
    \end{itemize}
\end{frame}

\begin{frame}
    \frametitle{Функции как аргументы}
    \begin{itemize}
        \item Пусть $A$ "--- любое множество. Определим
        \[T:(A\to\R)\to(A\to\R),\qquad T(f)=(x\mapsto f(x)+1).\]
        \item Что здесь является аргументом и значением?
        \item<2-> Аргумент $T$ "--- функция $f:A\to\R$. Значение $T(f)$ "--- другая функция $A\to\R$.
        \item<3-> Её значение на $x\in A$: $T(f)(x)=f(x)+1$.
        \item<4-> Полная лямбда-запись: $T=(f\mapsto(x\mapsto f(x)+1))$.
        \item<5-> Чтобы доказать равенство двух функций с одинаковыми доменом и кодоменом, достаточно проверить их значения на каждом аргументе.
    \end{itemize}
\end{frame}

\begin{frame}
    \frametitle{Свойства преобразования $T$}
    \begin{itemize}
        \item $T(f)=(x\mapsto f(x)+1)$. Проверим инъективность и сюръективность через обратную функцию.
        \item<2-> Вычитаем $1$ из каждого значения:
        \[\uncover<2->{S:(A\to\R)\to(A\to\R),\qquad S=(g\mapsto(x\mapsto g(x)-1)).}\]
        \item<3-> Для произвольных $f:A\to\R$ и $x\in A$:
        \[\uncover<4->{S(T(f))(x)=T(f)(x)-1=f(x)+1-1=f(x).}\]
        \item<5-> Для произвольных $g:A\to\R$, $x\in A$: $T(S(g))(x)=g(x)-1+1=g(x)$.
        \item<6-> Поэтому $S\circ T=\id_{A\to\R}$ и $T\circ S=\id_{A\to\R}$.
        \item<7-> Значит, $T$ "--- биекция и $T^{-1}=S$. При $A=\varnothing$ рассуждение тоже верно.
    \end{itemize}
\end{frame}

\begin{frame}
    \frametitle{Самостоятельная проверка}
    \begin{itemize}
        \item Определите, какие функции являются вложениями, наложениями, биекциями. Обоснуйте ответы.
        \item $u:\N\to\N$, $u(n)=n+1$.
        \item<2-> Вложение, но не наложение: $0$ не имеет прообраза.
        \item<3-> $v:\Z\to\Z$, $v(n)=|n|$.
        \item<4-> Не вложение ($v(1)=v(-1)$) и не наложение ($-1$ не имеет прообраза).
        \item<5-> $w:\R^2\to\R^2$, $w(x,y)=(x+y,y)$.
        \item<6-> Биекция. Обратная: $(a,b)\mapsto(a-b,b)$. Проверьте обе композиции.
        \item<7-> $f:\R\to\R$, $g:\Q\to\Q$, $f(x)=g(x)=x^3$.
        \item<8-> $f$ "--- биекция: куб строго возрастает, и для любого $b\in\R$ прообраз $\sqrt[3]{b}$ принадлежит $\R$.
        \item<9-> $g$ "--- вложение, но не наложение: у $2$ нет рационального прообраза ($\sqrt[3]{2}\notin\Q$).
    \end{itemize}
\end{frame}

\begin{frame}
    \frametitle{Итоги}
    \begin{itemize}
        \item Сначала проверяем, что правило задаёт функцию с указанными доменом и кодоменом.
        \item Значение не должно зависеть от представления аргумента.
        \item Инъективность доказываем для произвольных аргументов, опровергаем одной парой различных аргументов.
        \item Сюръективность доказываем построением прообраза произвольного элемента кодомена, опровергаем одним пропущенным элементом.
        \item Для обратной функции проверяем обе композиции.
        \item Для произвольных множеств не забываем пустые случаи.
    \end{itemize}
\end{frame}

\end{document}
